\begin{figure}[tb]
\centering
\includegraphics[width=.77\hsize]{images/chi_vertical.pdf}
\caption{
Cross-head importance (CHI).
In the above figures, ${}^th_i^l$ and $w_{ij}^l$ represents the output of the $i$-th neuron in the $l$-th layer just after training task $t$ and the parameter in $l$-th layer connecting between the neurons ${}^th_i^l$ to ${}^th_j^{l+1}$, respectively.
\textbf{(a)} The importance of $w_{ij}^l$ to task $\tau$ is computed based on its gradient, ${\partial \mathcal{L}^{\tau,\tau}}/{\partial w_{ij}^{l}}$ and then accumulated.
\textbf{(b)} After learning task $t$ ($t > \tau$), the state of related parameters might have been changed. To reflect importance to task $\tau$ again with the current neurons' output (e.g., ${}^th_i^l$ rather than old ${}^\tau h_i^l$), an additional loss, $\mathcal{L}^{t,\tau}$, is computed at task $\tau$'s head using task $t$'s data as unlabeled data for task $\tau$.
}
\label{fig:chi}
\vspace{-5mm}
\end{figure}